\section{Конструкторский раздел}
В данном разделе формулируются требования к разрабатываемому методу оценки кредитоспособности заемщика. Приводится описание метода, основанного на алгоритме градиентного бустинга деревьев решений, а также определяются его основные характеристики и область применения в задаче кредитного скоринга.

Рассматриваются этапы работы метода, включающие предобработку исходных данных, обучение модели и получение предсказаний кредитного риска. Описывается алгоритм реализации метода, включая последовательность вычислительных процедур и используемые математические модели. Раскрываются принципы интерпретации результатов моделирования.

Приводится описание структур данных, применяемых в системе, а также их назначение в процессе обработки информации. Определяется структура программного комплекса и описывается взаимодействие его компонентов, обеспечивающих реализацию разработанного метода.

%--------------------------
\subsection{Требования к разрабатываемому методу}

Метод оценки кредитоспособности заемщика должен удовлетворять следующим требованиям:

\begin{enumerate}
  \item Принимать на вход данные о заемщике, представленные в виде набора признаков.
  \item Обеспечивать предварительную обработку исходных данных, включающую преобразование признаков и обработку пропущенных значений.
  \item Осуществлять обучение модели на основе алгоритма градиентного бустинга деревьев решений.
  \item Формировать предсказание кредитоспособности заемщика в виде бинарной классификации.
  \item Обеспечивать возможность интерпретации результатов работы модели, включая определение значимости признаков.
  \item Обеспечивать достаточную точность прогнозирования для применения в задачах кредитного скоринга.
\end{enumerate}

%-------------------------
\subsection{Требования к программному комплексу}

Программный комплекс, реализующий интерфейс к разработанному методу оценки кредитоспособности заемщика, должен обеспечивать выполнение следующих функций:

\begin{itemize}
  \item загрузку исходных данных о заемщиках с использованием пользовательского интерфейса;
  \item выполнение предварительной обработки данных, включающей преобразование признаков и подготовку данных к обучению модели;
  \item запуск процесса обучения модели на основе алгоритма градиентного бустинга деревьев решений;
  \item выполнение предсказания кредитоспособности заемщика на основе обученной модели;
  \item отображение результатов работы метода, включая значение скорингового балла и оценку кредитного риска;
  \item отображение значимости признаков, используемых моделью, для интерпретации результатов.
\end{itemize}

Входные данные загружаются из файловой системы в формате JSON. Система поддерживает повторный запуск скоринга после изменения входных данных или параметров модели. При повторном запуске формируется новый входной JSON-файл, который передается в модуль предсказания.

%--------------------------
\subsection{Основной алгоритм работы}

Учитывая особенности существующих методов оценки кредитоспособности, рассмотренные ранее, был разработан метод, основанный на алгоритме градиентного бустинга деревьев решений.

На первом этапе осуществляется загрузка и подготовка исходных данных о заемщиках. Данные приводятся к требуемому формату, выполняется обработка пропущенных значений, преобразование признаков и разделение данных на обучающую и тестовую выборки, необходимых для последующего обучения модели.

Далее производится обучение модели на основе подготовленных данных. В процессе обучения формируется ансамбль решающих деревьев, позволяющий учитывать сложные зависимости между признаками и повышать точность прогнозирования. Обучение продолжается до достижения заданного критерия качества модели.

После завершения обучения модель сохраняется и в дальнейшем используется для формирования предсказаний. Для этого выполняется загрузка обученной модели и обработка новых данных о заемщике, прошедших этап предобработки.

Модель по входным данным выполняет бинарную классификацию и относит заемщика к одному из классов: кредитоспособен или некредитоспособен. Результат работы модели может быть дополнительно проанализирован с использованием методов интерпретации, включая оценку значимости признаков. Такой анализ позволяет определить, какие характеристики заемщика оказывают наибольшее влияние на итоговое предсказание модели, и повышает прозрачность применения алгоритмов машинного обучения в задаче кредитного скоринга \cite{MolnarInterpretability}.

Схема алгоритма представлена на рисунке \ref{contruct::gbdt_credit_scoring_algorithm}.


\FloatBarrier
\begin{figure}[!ht]
  \begin{center}
    \includegraphics[width=\linewidth, height=0.5\textheight, keepaspectratio]{inc/png/gbdt_credit_scoring_algorithm.png.drawio.png}
  \end{center}
  \captionsetup{justification=centering}
  \caption{Верхнеуровневая схема алгоритма}
  \label{contruct::gbdt_credit_scoring_algorithm}
\end{figure}
\FloatBarrier


\subsection{Описание входных данных и признакового пространства}

Входные данные для разработанного метода оценки кредитоспособности представляют собой совокупность характеристик физического лица, отражающих его финансовое состояние, социально-демографические особенности и кредитную историю.

Каждое физическое лицо описывается вектором признаков следующего вида:
\begin{equation}
  x = (x_1, x_2, \ldots, x_d),
\end{equation}

где

\begin{itemize}
  \item $x$ --- вектор признаков объекта;
  \item $x_i$ --- значение $i$-го признака;
  \item $d$ --- общее количество признаков.
\end{itemize}

В рамках разработанного метода используется набор признаков, включающий следующие группы:

\begin{itemize}
  \item \textbf{Социально-демографические характеристики:} возраст, семейное положение, уровень образования и другие параметры заемщика;

  \item \textbf{Финансовые характеристики:} уровень дохода, параметры запрашиваемого кредита, размер платежей и долговая нагрузка;

  \item \textbf{Параметры занятости:} тип занятости, стаж работы и характеристики источника дохода;

  \item \textbf{Кредитная история:} сведения о ранее полученных кредитах, наличии задолженностей и просрочек;

  \item \textbf{Агрегированные поведенческие признаки:} показатели, полученные на основе истории предыдущих заявок и фактического выполнения платежей по кредитам.
\end{itemize}

Признаки характеризуются различным типом данных: они могут быть количественными (например, уровень дохода, величина расходов) и категориальными (например, тип занятости, уровень образования). В связи с этим на этапе предобработки требуется их дополнительное преобразование.

Целевая переменная задачи определяется как бинарный признак:
\begin{equation}
  y \in \{0, 1\},
\end{equation}

где

\begin{itemize}
  \item $y$ --- метка класса;
  \item $y = 1$ --- заемщик с низким риском дефолта (кредитоспособный);
  \item $y = 0$ --- заемщик с высоким риском дефолта (некредитоспособный).
\end{itemize}


Таким образом, задача оценки кредитоспособности формализуется как задача бинарной классификации, в которой по вектору признаков $x$ необходимо определить значение целевой переменной $y$.


\subsection{Предобработка данных}

Для обеспечения корректной работы модели и повышения качества прогнозирования выполняется этап предобработки исходных данных. В задачах прогнозного моделирования предобработка является одним из ключевых этапов, поскольку исходные данные могут содержать пропущенные значения, выбросы, признаки различного масштаба и категориальные переменные. Их предварительное преобразование позволяет подготовить данные к обучению модели и снизить влияние некорректных или неполных наблюдений на итоговое качество прогнозирования \cite{KuhnPreprocessing}.

В разработанном методе предобработка состоит из следующих этапов:

\begin{itemize}
  \item обработка пропущенных значений предполагает заполнение пропусков медианными значениями для числовых признаков и часто встречающимися значениями для категориальных признаков, что позволяет снизить влияние неполноты данных на качество модели;

  \item обработка выбросов выполняется путем ограничения значений признаков по заданным порогам (квантильное усечение), что повышает устойчивость модели к экстремальным значениям доходов и других финансовых характеристик заемщика;

  \item кодирование категориальных признаков осуществляется с использованием метода one-hot кодирования, что обеспечивает корректную обработку категориальных данных алгоритмом машинного обучения;

  \item масштабирование признаков при необходимости используется для приведения числовых признаков к сопоставимому диапазону значений, однако для алгоритма градиентного бустинга данный этап не является обязательным;

  \item формирование производных признаков включает расчет дополнительных показателей, отражающих финансовое состояние заемщика, таких как:
        \begin{itemize}
          \item отношение ежемесячных платежей по кредитам к доходу (DTI);
          \item отношение расходов к доходу;
          \item коэффициент долговой нагрузки;
        \end{itemize}
\end{itemize}

После выполнения предобработки формируется итоговый набор признаков, используемый для обучения модели и построения прогнозов кредитоспособности.


\subsection{Параметры и настройка модели}

В рамках разработанного метода для оценки кредитоспособности используется модель градиентного бустинга деревьев решений. Качество такой модели во многом определяется выбором гиперпараметров, включая количество деревьев, глубину деревьев, скорость обучения и параметры регуляризации. Эти параметры влияют на сложность ансамбля, скорость обучения и способность модели обобщать закономерности на новых данных \cite{HastieElementsBoosting}.

К основным параметрам модели относятся:

\begin{itemize}
  \item количество деревьев $M$ определяет размер ансамбля; увеличение числа деревьев позволяет повысить точность модели, однако приводит к увеличению времени обучения и риску переобучения;

  \item глубина деревьев ограничивает сложность базовых алгоритмов; небольшая глубина способствует лучшей обобщающей способности модели, тогда как чрезмерная глубина может привести к переобучению;

  \item скорость обучения $\gamma$ определяет вклад каждого дерева в итоговую модель; малые значения обеспечивают более стабильное обучение, но требуют увеличения числа итераций;

  \item минимальное количество объектов в листе дает возможность контролировать процесс разбиения деревьев и снижает влияние шумовых данных;

  \item параметры регуляризации используются для предотвращения переобучения модели и повышения ее обобщающей способности;
\end{itemize}

Настройка гиперпараметров осуществляется на основе обучающей выборки с использованием процедуры валидации. Разделение исходных данных на обучающую и тестовую выборки позволяет сначала построить модель на одной части данных, а затем оценить ее качество на независимых наблюдениях, которые не использовались при обучении \cite{BishopPatternRecognition}.

В процессе настройки подбираются такие значения параметров, при которых достигается наилучшее качество классификации при сохранении устойчивости модели к переобучению.

\subsection{Формирование результата и принятие решения}

После завершения обучения модель градиентного бустинга используется для оценки кредитоспособности физического лица на основе входных данных.

На вход модели подается вектор признаков заемщика $x$, прошедший этап предобработки. В результате модель формирует числовую оценку, которая интерпретируется как вероятность наступления события дефолта:
\begin{equation}
  P(y = 1 \mid x),
\end{equation}

где $P(y = 1 \mid x)$ --- вероятность того, что заемщик с признаками $x$ относится к классу кредитоспособных заемщиков.

Для удобства интерпретации полученная вероятность может быть преобразована в скоринговый балл, отражающий уровень кредитного риска заемщика.

Принятие решения о предоставлении кредита осуществляется на основе порогового значения вероятности. Введем порог $T$, определяемый кредитной политикой банка.

Решение принимается следующим образом:
\begin{equation}
  \hat{y} =
  \begin{cases}
    1, & \text{если } P(y = 1 \mid x) \geq T, \\
    0, & \text{если } P(y = 1 \mid x) < T.
  \end{cases}
\end{equation}

Таким образом, если вероятность кредитоспособности превышает установленный порог, заемщик относится к категории надежных и принимается положительное решение о выдаче кредита. В противном случае заемщик классифицируется как рискованный, и в выдаче кредита может быть отказано.

Пороговое значение $T$ может варьироваться в зависимости от кредитной политики банка, уровня допустимого риска и текущей экономической ситуации.

Дополнительно результат работы модели может использоваться для ранжирования заемщиков по уровню риска и определения индивидуальных условий кредитования, таких как процентная ставка или размер кредитного лимита.


\subsection{Структуры данных}

При реализации метода оценки кредитоспособности применяются структуры данных для хранения, обработки и передачи информации между компонентами системы.

Основной структурой является набор данных о физических лицах, представленный в виде таблицы, где каждая строка соответствует отдельному заемщику, а столбцы — признакам, характеризующим его финансовое состояние и поведенческие характеристики.

Для представления входных данных используется следующая структура:
\begin{itemize}
  \item идентификатор заемщика;
  \item значения признаков (социально-демографические параметры, кредитная история, финансовые характеристики а также агрегированные поведенческие показатели заемщика);
  \item целевая переменная (для обучающей выборки).
\end{itemize}

Обучающая выборка формируется в виде матрицы признаков $X$ и вектора целевой переменной $y$:
\begin{equation}
  X \in \mathbb{R}^{n \times d}, \quad y \in \{0, 1\}^{n},
\end{equation}

где

\begin{itemize}
  \item $X$ --- матрица признаков;
  \item $y$ --- вектор целевых значений;
  \item $n$ --- количество объектов;
\end{itemize}

Модель градиентного бустинга представляется в виде ансамбля решающих деревьев, параметры которых включают структуру деревьев, значения весов и коэффициенты модели.

Результат работы модели представляет собой структуру данных, включающую вероятность кредитоспособности, скоринговый балл, класс заемщика, а также значения важности признаков.

Для хранения промежуточных данных используются структуры, содержащие результаты предобработки, а также параметры модели, полученные в процессе обучения.

\subsection{Взаимодействие компонентов системы}

Метод реализован как программный комплекс из взаимосвязанных компонентов, выполняющих полный цикл обработки данных и принятия решений. В структуре системы выделяются следующие основные модули:

\begin{itemize}
  \item Модуль загрузки данных обеспечивает получение исходной информации о физических лицах из внешних источников или пользовательского интерфейса.

  \item Модуль предобработки данных выполняет очистку данных, кодирование признаков, обработку пропущенных значений и формирование итогового набора признаков.

  \item Модуль обучения модели реализует процесс построения модели на основе алгоритма градиентного бустинга деревьев решений с использованием обучающей выборки.

  \item Модуль предсказания обеспечивает применение обученной модели к новым данным и расчет вероятности кредитоспособности заемщика.

  \item Модуль интерпретации результатов выполняет анализ важности признаков и формирование объяснений полученных результатов.

  \item Модуль визуализации и вывода результатов отображает результаты работы системы, включая скоринговый балл, оценку кредитного риска и рекомендации по принятию решения.
\end{itemize}

Система работает в двух режимах: режиме предсказания (online) и режиме обучения модели (offline). В режиме обучения используются исторические данные о физических лицах. Эти данные проходят этап предобработки и поступают в модуль обучения, где формируется модель градиентного бустинга. В режиме предсказания новые данные о заемщике проходят этап предобработки и передаются в модуль предсказания, который использует ранее обученную модель для расчета вероятности кредитоспособности.

Взаимодействие компонентов осуществляется последовательно. Исходные данные поступают в модуль загрузки, после чего передаются в модуль предобработки. Подготовленные данные используются либо для обучения модели, либо для получения предсказания. Результаты работы модели поступают в модуль интерпретации, где определяется значимость признаков, а затем передаются в модуль визуализации для отображения пользователю.

На рисунке \ref{fig:system_architecture} представлена схема взаимодействия компонентов системы.


\FloatBarrier
\begin{figure}[!ht]
  \begin{center}
    \includegraphics[width=\linewidth, height=0.4\textheight, keepaspectratio]{inc/pdf/Схема взаимодействия компонентов системы.pdf}
  \end{center}
  \captionsetup{justification=centering}
  \caption{Схема взаимодействия компонентов системы}
  \label{fig:system_architecture}
\end{figure}
\FloatBarrier

\subsection*{Выводы}
В данном разделе был разработан метод оценки кредитоспособности физических лиц на основе алгоритма градиентного бустинга деревьев решений. Были сформулированы требования к методу и программному комплексу, а также описаны основные этапы его функционирования.

Представлено описание признакового пространства, включающего финансовые характеристики, социально-демографические параметры, кредитную историю, а также поведенческие признаки заемщиков. Рассмотрены методы предобработки данных, направленные на повышение качества модели. Рассмотрены этапы обучения модели, формирования предсказаний и принятия решения на основе порога вероятности.

Также были описаны используемые структуры данных и разработана архитектура программного комплекса, включающая взаимодействие его основных компонентов. Для наглядного представления работы системы приведена структурная схема.
